\chapter{Example: the high-frequency limit}\label{ch:ex-highfreq}

When the drive is faster than every internal process, the system responds to a
time-averaged, renormalized Hamiltonian. This limit is the basis of Floquet
engineering~\cite{goldman2014periodically,bukov2015universal,eckardt2017colloquium}, and it is
a direct test of the van Vleck expansion of \cref{sec:van-vleck}.

\section{Accuracy of the van Vleck expansion}

\Cref{fig:highfreq}(a) compares the eigenvalues of the van Vleck Hamiltonians of order
$n=0$, $1$ and $2$ with exact Sambe quasienergies for a single-lead model driven
simultaneously through its phase and level energy, with a phase offset between the two
drives. The error decreases as $\Omega^{-(n+1)}$, as the expansion
predicts~\cite{eckardt2015high,mikami2016brillouin}. With a relative phase the drive is not
even in time, $H_{-k}\ne H_k$, and the first-order term $[H_1,H_{-1}]/\Omega$ is finite. For a
single-parameter cosine drive it vanishes and the zeroth-order result is already accurate to
$O(\Omega^{-2})$.

\section{Renormalized coupling and coherent destruction of tunnelling}

For a phase drive with $\Omega$ larger than the many-body spectral width, the zeroth-order
Hamiltonian is $H_0$ of \cref{eq:phase-h0}: the junction with
\begin{equation}\label{eq:gamma-renorm}
  \Gamma\to\Gamma_{\mathrm{eff}}=\Gamma J_0\!\left(\frac A4\right)^2.
\end{equation}
Each tunnelling amplitude carries the phase factor $\ee^{\mp\ii\phi(t)/4}$. Its time average is
$J_0(A/4)\,\ee^{\mp\ii\phi_0/4}$, so the hybridization, quadratic in the tunnelling amplitude,
is reduced by $J_0(A/4)^2$. As $A$ increases, the dot decouples progressively from the leads; at
the first zero of $J_0$, $A=4\times2.405\approx9.62$, it is isolated to leading order in
$1/\Omega$. This is the junction version of dynamical
localization~\cite{dunlap1986dynamic,holthaus1992collapse,lignier2007dynamical} and of
coherent destruction of tunnelling~\cite{grossmann1991coherent,kayanuma1994role}.

\Cref{fig:highfreq}(b) follows the quasienergies of the lowest singlets and of the doublet
for $U=\q{highfreq.u}$, $\Gamma=\q{highfreq.gamma}$, $\eps=\q{highfreq.detuning}$ and
$\Omega=\q{highfreq.omega}$, together with the static energies of the junction with the
renormalized coupling \eqref{eq:gamma-renorm}. The undriven ground state is a singlet. As
$\Gamma_{\mathrm{eff}}$ decreases, the doublet becomes the lowest dressed state near
$A\approx\q{highfreq.crossing}$: the drive induces a level crossing between the dressed
singlet and doublet, a dynamical analogue of the 0--$\pi$ transition. At the zero of $J_0$
the time-averaged coupling vanishes and the atomic energies are recovered to leading order,
$\qe_{S,S^*}=U/2$ and $\qe_D=0$. The remaining small deviations, which also split $S$ and $S^*$,
are the $O(\Omega^{-2})$ corrections of the higher harmonics $J_k(A/4)$, $k\ge1$; the
second-order van Vleck Hamiltonian reproduces them.

\begin{figure}[htb]
\centering
\manualfigure{highfreq}
\caption{High-frequency limit. (a) Maximum deviation between the eigenvalues of the van
Vleck Hamiltonian of order $n$ and exact Sambe quasienergies, for a single-lead model with
combined phase and level drive; dotted lines are $\Omega^{-(n+1)}$. (b) Quasienergies of the
phase-driven junction ($U=\q{highfreq.u}$, $\Gamma=\q{highfreq.gamma}$,
$\eps=\q{highfreq.detuning}$, $\phi_0=\q{highfreq.phi}$, $\q{highfreq.pairs}$ bath level pair
per lead with $D=\q{highfreq.bandwidth}$, $\Omega=\q{highfreq.omega}$) versus drive amplitude.
Markers are Floquet results, lines the static energies with $\Gamma\to\Gamma J_0(A/4)^2$.}
\label{fig:highfreq}
\end{figure}

\begin{caution}
The renormalization \eqref{eq:gamma-renorm} describes the formal limit with $\Omega$
above the full reservoir bandwidth. Microwave experiments on superconducting
junctions typically operate below that bandwidth. For
$\Gamma,\abs\eps\ll\Omega\ll\Delta$ the leads follow the drive adiabatically, and the
relevant time average is that of the induced pairing $\Gamma_\phi(t)$,
$\Gamma\cos(\phi_0/2)\to\Gamma J_0(A/2)\cos(\phi_0/2)$ (\cref{sec:driven-largegap}). The
order of limits matters, and the van Vleck expansion of the full model is a reliable
approximation only in the first regime.
\end{caution}

\begin{incode}
\begin{lstlisting}[frame=none,backgroundcolor={},aboveskip=0pt,belowskip=0pt]
fast = dict(bath=cosh_grid(pairs=1, bandwidth=2.), u=1., gamma=.8, phi=.8, detuning=0.,
            symmetry=False, compress=False)
periodic = PeriodicHamiltonian.from_parameters(reference_model, fast,
                                               ParameterDrive.cosine("phi", 6.), omega=40.)
effective = periodic.van_vleck(2)                     # a static Hamiltonian
energies = solve(effective, 2, Sector(1, 1)).energies    # dressed doublet, panel (b) of the figure
\end{lstlisting}
The effective Hamiltonian keeps the coordinates of the model, so it can be compared with
static calculations and solved by DMRG.
\par\smallskip\usermanual{docs/floquet.md}.
\end{incode}

\tutorial{04_high_frequency}
